\chapter{绪论}
\section{研究背景与意义}
自习室为学生提供了相对安静的学习空间。考试周等集中使用时段中，学生往往需要到现场才能确认座位是否可用。若预约信息仅依靠人工登记维护，同一座位的时间安排容易出现遗漏，取消预约后也可能仍被标记为占用。将座位与时间区间统一管理，有助于减少重复登记，并使查询结果能够反映当前预约状态。

软件工程强调从需求出发组织设计、实现与验证活动\cite{sommerville2016}。本文以小规模预约原型为研究对象，将业务规则转化为可执行的校验条件，通过固定输入与预期输出检验核心行为。研究重点是时段冲突、取消权限和数据一致性，原型采用本地命令行入口调用业务函数。

\section{研究内容与方法}
\subsection{研究内容}
本文完成三方面工作：一是明确学生预约座位时需要遵守的时间与归属规则；二是建立预约数据模型，实现创建、取消和查询功能；三是设计可重复执行的功能用例，核对预约结果与数据库记录。业务原型与测试程序随论文源码保存，便于重新运行和检查。

\subsection{研究方法}
需求分析采用典型使用场景，设计阶段采用分层结构与关系数据模型，实现阶段使用参数化 SQL 访问数据库，验证阶段采用功能用例逐项核对结果。测试数据在每次运行时重新初始化，以减少历史记录对结果的影响。

\section{论文组织结构}
第二章介绍业务需求和约束，第三章给出系统结构、数据模型与冲突判定方法，第四章说明预约和取消的实现，第五章报告测试与统计结果，第六章总结工作并提出后续改进方向。附录列出运行方法及核心源代码。
